\begin{claimboundary}
Realtimeという言葉はlatencyの数字を強く印象付けるが、Mistralが示したsub-200ms latency、WER、price/performance比較はvendor-reportedであり、記載されたserving configurationに依存する。本稿ではopen-weight・license・model size・language scopeといったrelease factと、速度・精度の優位性を同一視しない。またopen weightsであることは、周辺のhosted service全体がopen sourceであることを意味しない。一次資料はMistralのnews indexを起点としているため、item-levelの表現もFebruary launch articleの範囲から広げない。\autocite{src-d8b7c1c36830}
\end{claimboundary}

Google側では、2月19日にGemini 3.1 Pro Previewとcustom-tools endpoint、2月26日にGemini 3.1 Flash Image Preview（Nano Banana 2）がGemini API changelogへ記録された。前者はreasoning/tool-use側、後者はspeed/high-volume image generation側の更新であり、一つの『Gemini 3.1 update』として潰すべきではない。特に本章の主役はinteractionのmodal breadthなので、Pro Previewはtool surfaceの変化を示すcontext、Flash Image Previewは画像生成の応答面を広げる主題として役割を分ける。\autocite{src-a24fd97eb3d0}


\begin{claimboundary}
Gemini 3.1 ProとGemini 3.1 Flash Imageは、2月時点ではいずれもPreviewとして記録されている。したがって、後日のstable availabilityを2月へ遡及させたり、changelogの存在だけを根拠にquality・benchmark superiorityを独立検証済みとみなしたりはしない。API changelogが強いのはidentifierとchronologyの固定であり、性能評価の再現ではない。\autocite{src-a24fd97eb3d0}
\end{claimboundary}

model releaseがcommon implementation layerへ到達したことも確認できる。Transformers v5.2.0は2月16日にVoxtralRealtime、GLM-5、Qwen3.5/Qwen3.5 MoE、VibeVoice Acoustic Tokenizerをsupportへ加え、VoxtralRealtimeについてはaudio encoderとMistral-based LM decoder、time conditioning、causal-convolution cachesを使うstreaming ASR implementationを説明している。これはMistralのlatencyやaccuracy claimを再検証するものではないが、realtime speechが一般的なframework側の実装対象になったことを示す。\autocite{src-0bd04783f97b,src-a43d64987293}


\begin{claimboundary}
Transformers integrationは『実装可能になった』ことの証拠であって、model vendorのbenchmarkを独立に保証するものではない。v5.1.0は同月前半のsupporting context、v5.2.0がより強いFebruary milestoneである。またv5.2にはbreaking interface changesも含まれるため、新model supportだけを切り出して互換性が全面的に単純化したと読むこともできない。\autocite{src-0bd04783f97b,src-a43d64987293}
\end{claimboundary}

2月のinteraction進化をまとめると、二つの方向が見える。Voxtral Realtimeは音声をstreamとして処理し、Gemini 3.1 Flash Image Previewはimage generationを高速・大量利用向けのAPI surfaceへ置く。これは『より賢いtext model』とは別の競争であり、input/outputのmodalityと応答時間をworkloadに合わせて組み替える動きである。3月以降を読む際も、model intelligenceだけでなく、音声・画像・toolを含むinteraction loopがどのruntimeとAPIへ定着するかが独立した観察軸になる。\autocite{src-a24fd97eb3d0,src-a43d64987293,src-d8b7c1c36830}
